\documentclass[10pt,a4paper,twocolumn]{article}
\usepackage[T1]{fontenc}
\usepackage{lmodern}
\usepackage{amsmath}
\usepackage{graphicx}
\usepackage{booktabs}
\usepackage[margin=1.8cm,columnsep=0.7cm]{geometry}
\usepackage[colorlinks=true,linkcolor=black,citecolor=black,urlcolor=blue]{hyperref}
\usepackage{caption}
\usepackage{titlesec}
\usepackage{microtype}
\captionsetup{font=small,labelfont=bf,labelsep=period}
\titleformat{\section}{\normalfont\bfseries\MakeUppercase}{\thesection}{0.6em}{}
\titleformat{\subsection}{\normalfont\itshape}{\thesubsection}{0.6em}{}
\titleformat{\subsubsection}{\normalfont\itshape\small}{\thesubsubsection}{0.6em}{}
\setlength{\parskip}{0pt}
\renewcommand{\arraystretch}{1.05}

\begin{document}
\twocolumn[{
\begin{center}
{\small Draft report, 30 October 2026 \hfill ComPy 2.4.0}\\[10pt]
{\LARGE\bfseries Shallow crustal shear velocity of $\sim$40\,Ma Pacific lithosphere\\[3pt]
from seafloor compliance at the Young ORCA array}\\[12pt]
{\large M.\,A. Aminian}\\[2pt]
{\small Institut de Physique du Globe de Paris, Universit\'e Paris Cit\'e. Processing, figures and first draft
prepared with Claude Code from the ComPy pipeline; to be checked by the author.}\\[12pt]
\begin{minipage}{0.86\textwidth}
\small
\textbf{SUMMARY}\\
We apply the compliance method of Aminian et al.\ (2025), now implemented end to end in ComPy 2.4.0,
to a second broad-band ocean-bottom network: the 30-station Young ORCA array (FDSN code XE, 2018--2019)
on $\sim$40\,Ma Pacific seafloor formed at the fast-spreading Pacific--Farallon ridge, at water depths of
4386--4824\,m. The stations carry Nanometrics Trillium 240 seismometers and differential pressure gauges,
as do the RHUM-RUM stations of the earlier study. After removing clipped and dead data, earthquakes, tilt and
horizontal-coherent noise, and calibrating each pressure gauge on Rayleigh waves of $M\geq7$ earthquakes
(gain factors 0.70--1.24, $\pm$1.4\,per\,cent between events), we measure the normalized compliance
between 4.9 and 14.0\,mHz at 20 stations from 5038 35-min windows per station
(median), with coherence above 0.9 and month-to-month variations of a few per cent. A Bayesian inversion
(delayed-rejection adaptive Metropolis with parallel tempering, free crustal Vp/Vs, pressure gain sampled within
$\pm$5\,per\,cent) fits every station (median $\chi^2/N$ = 0.36). Below 79--210\,m of sediment,
consistent with CRUST1.0, the network-pooled shear velocity rises from 2.37\,km\,s$^{-1}$ at 0.3\,km to
2.94\,km\,s$^{-1}$ at 1\,km and 3.23\,km\,s$^{-1}$ at 2\,km below the seafloor, slower than CRUST1.0
in the upper crust. Lower-crustal Vs depends on the Vp/Vs assumption (3.54 with Vp/Vs free,
4.05\,km\,s$^{-1}$ with Vp/Vs = 1.85). Converted with the relation used for the Indian Ocean ridges,
the upper-crustal velocities would imply 63 (27--91)\,per\,cent serpentinization at 1\,km; on magmatic crust formed at a fast-spreading
ridge they instead record the porosity of the extrusive section. This normal-crust reference shows that
compliance-derived Vs alone cannot identify serpentinite, and it calibrates the interpretation of the
low velocities found at the South-West Indian Ridge.\\[4pt]
\textbf{Key words:} Oceanic crust; Seafloor compliance; Inversion; Pacific Ocean; Monte Carlo methods.
\end{minipage}
\end{center}
\vspace{14pt}
}]

\section{Introduction}
Seafloor compliance, the transfer function between vertical seafloor displacement and the pressure of
infragravity waves, constrains the shear velocity (Vs) of the upper few kilometres beneath a single broad-band
ocean-bottom seismometer (OBS) equipped with a differential pressure gauge (DPG) (Crawford et al.\ 1991, 1998).
Aminian et al.\ (2025) measured it at eight RHUM-RUM stations in the Indian Ocean and interpreted very low upper-crustal Vs
beneath the ultraslow South-West Indian Ridge (SWIR) as serpentinized mantle-derived peridotite.

Two questions follow. Is the processing reproducible on another network without manual tuning? And what does
the same analysis give on crust that is \emph{not} expected to contain serpentinite? The Young ORCA experiment
(Eilon et al.\ 2022) answers both. Its 30 broad-band OBS were deployed for about 13 months in 2018--2019 on
$\sim$40\,Ma Pacific lithosphere, formed at the fast-spreading Pacific--Farallon ridge, a setting where the crust is
fully magmatic and layered in the classical way (e.g.\ Christeson et al.\ 2019).

This report applies ComPy 2.4.0, which now runs from the raw archive to compliance, Vs and
serpentinization unattended, to the whole network, using the method of Aminian et al.\ (2025) with the corrections
of ComPy 2.2--2.4 (Section~4). We present the compliance functions, the Vs models and their resolution, and
discuss what they imply for the serpentinization interpretation.

\section{Theory}
The vertical compliance $\zeta(\omega)$ is the ratio of the spectra of vertical displacement and vertical
stress at the seafloor (Sorrells \& Goforth 1973). We use the normalized compliance $\eta(\omega) = k(\omega)\zeta(\omega)$,
with $k$ the infragravity wavenumber from $\omega^2 = gk\tanh(kH)$ for water depth $H$. The measured compliance is
\begin{equation}
\hat\eta(\omega) = k\,\frac{\omega^2\,|\hat S_{pz}|/\hat S_{pp} + a_w/p}{\omega^2 + 2g/R},
\end{equation}
where $\hat S_{pz}$ and $\hat S_{pp}$ are the cross- and pressure spectra stacked over the selected windows,
$a_w/p = 2\pi (G/g)\,e^{-kH}\cosh kH$ is the gravitational attraction of the passing water mass per unit pressure, and $2g/R$
the free-air gradient (Crawford et al.\ 1998). The ratio $|\hat S_{pz}|/\hat S_{pp}$ equals the coherence times
$\sqrt{\hat S_{zz}/\hat S_{pp}}$ used in the paper's eq.~(6). The wave-attraction term is deterministic and is not
multiplied by the coherence: earlier ComPy versions did so (2.0--2.1) or dropped it (before 2.0), which biased $\eta$
low at the lowest frequencies. The random error follows Bendat \& Piersol (1993),
\begin{equation}
\epsilon[\hat\eta] = \frac{\sqrt{1-\hat\gamma^2}}{|\hat\gamma|\sqrt{2n_d}}\,|\hat\eta|,
\end{equation}
with $n_d$ the number of independent spectral averages over all accepted windows.

The forward model is the minor-vector propagator of Crawford et al.\ (1991) for a layered elastic half-space.
The likelihood is Gaussian with a full data covariance $\mathbf{C}$,
\begin{equation}
L(\mathbf{m}) \propto \exp\!\left[-\tfrac12\,(\hat{\boldsymbol\eta} - \boldsymbol\eta(\mathbf{m}))^{\mathsf T}\mathbf{C}^{-1}(\hat{\boldsymbol\eta} - \boldsymbol\eta(\mathbf{m}))\right],
\end{equation}
where $\mathbf{C}$ sums the random error, the covariance of monthly stacks (which measures non-stationarity and
correlates neighbouring frequencies), and a 1\,per\,cent floor for 2-D structure. The pressure gain is a parameter
of the model, so the 5\,per\,cent calibration uncertainty is shared by all frequencies and is not counted once per frequency.

Serpentinization $\beta$ is converted from Vs with the linear relation of Miller \& Christensen (1997) used in
Aminian et al.\ (2025, Fig.~10):
\begin{equation}
\beta = 100\,\frac{4.2 - V_s}{4.2 - 2.2}\ \text{per cent}\quad (V_s \text{ in km\,s}^{-1}),
\end{equation}
clipped to 0--100 and evaluated on every posterior draw.

\section{Data set}
Young ORCA (Gaherty et al.\ 2019; Eilon et al.\ 2022) deployed 30 Scripps broad-band OBS between
4$^\circ$S and 8$^\circ$S, 130$^\circ$--136$^\circ$W (Fig.~\ref{fig:map}), in water 4386--4824\,m deep.
Each carries a Trillium 240 seismometer and a DPG, recorded at 50\,Hz. The data come from the BRUIT-FM archive
(SDS, 587\,GB for XE), with responses from the EarthScope/IRIS DMC. Because reading the archive disk in parallel was slow,
the raw files of most stations were first copied to local disk station by station; the processing is identical.

Three stations were not processed: CC02, whose pressure channel does not work, and CC10 and CC12, which recorded for
50 and 17 days. Of the 27 processed stations, 3 have no 13-h span in which all four channels are usable
(EC02, EE02, EE03; the seismometer is clipped or dead most of the time, Section~4.1), at WW02 no window passes
the quality gates, and 3 do not reach the 500 windows required by Aminian et al.\ (2025): EC03 (274), EC04 (5), WC01 (3). Table~\ref{tab:stations} lists the 20 stations with a compliance measurement; 20 of them are inverted.

\begin{figure}
\centering
\includegraphics[width=\columnwidth]{fig01_map.png}
\caption{Young ORCA (XE) stations on bathymetry (ETOPO1). Red: compliance measured and inverted; orange: compliance
measured but rejected or not inverted; grey: not processed or no usable data. Inset: location in the Pacific.}
\label{fig:map}
\end{figure}

\section{Data processing, compliance measurement and inversion}
The processing follows the flowchart of Aminian et al.\ (2025, Fig.~2), run unattended by \texttt{ComPy/pipeline}.

\subsection{Pre-processing}
\textit{Decimation, clipping and responses.} The raw data are decimated to 1.25\,Hz with FIR filters (factors 5, 4, 2).
Many Young ORCA seismometers sat at full scale ($2^{23}$ counts) for weeks, presumably while the sensor was off level, and
some channels went flat. Following the paper's rule that windows with a saturated component are discarded, every
10-min block in which any channel has a sample within 1\,per\,cent of full scale, or a standard deviation below 5 counts, is
cut from all channels, with 5 minutes on each side, before decimation. This removed 28\,per\,cent of the record at the
processed stations. Responses are removed after decimation without a water level and with an explicit pre-filter
(0.8--600\,mHz corners), giving displacement and pressure.

\textit{Earthquakes and local events.} Spans after $M\geq5.85$ earthquakes in the USGS catalogue (1.5\,d per magnitude unit;
Deen et al.\ 2017) and local events detected with a recursive STA/LTA trigger (0.1--0.5\,Hz) are excluded from the spectral
estimates.

\textit{Glitches.} The hourly mass-position glitch of the RHUM-RUM instruments is absent here: folding 15-day blocks at
trial periods around 1\,h gave fold strengths of 1.15 (median) and 1.76 (maximum) at six stations, below the 2.0 needed to
subtract a template, so the step was switched off for the rest of the network.

\textit{Tilt and coherent noise.} Hourly tilt angles that minimize the variance of the vertical (Harmon et al.\ 2022) are
median-smoothed over 24\,h and applied per sample, so the rotation adds no steps at hour boundaries (Fig.~\ref{fig:tilt}).
The remaining horizontal-coherent noise is then removed from the vertical with transfer functions estimated over
$\leq$24\,h windows (Crawford \& Webb 2000). Fig.~\ref{fig:psd} shows the vertical acceleration PSD at CC11 before
and after cleaning; the compliance bump near 0.01\,Hz stands at -158\,dB.

\begin{figure*}
\centering
\includegraphics[width=0.92\textwidth]{fig03_psd_stages.png}
\caption{(a) Vertical acceleration PSD at CC11 after response removal (blue) and after tilt and coherent-noise
removal (orange); median and 10--90 per cent range of 1-h windows; green: compliance band. (b) Calibrated pressure PSD.}
\label{fig:psd}
\end{figure*}

\begin{figure}
\centering
\includegraphics[width=\columnwidth]{fig04_tilt.png}
\caption{Hourly tilt azimuth (a) and incidence angle (b) at CC11, coloured by the logarithm of the variance reduction.}
\label{fig:tilt}
\end{figure}

\subsection{Pressure gauge calibration}
Each DPG is calibrated in situ on the Rayleigh waves of $M\geq7$ earthquakes (Zha \& Webb 2016; Aminian et al.\ 2025): the measured
pressure-to-acceleration ratio, normalized by $\rho_w H$, is fitted to its theoretical value between 0.02 and 0.06\,Hz for events
whose magnitude-squared coherence exceeds 0.97 (Fig.~\ref{fig:cal}). Gains were obtained at 18 of the 20 stations with a compliance measurement,
from 1--12 events each, and range from 0.70--1.24 with an event-to-event spread of 0.1--3.2\,per\,cent.
Where the BRUIT-FM microseism calibration exists, the ratio of the two gains is 0.986 (0.835--1.029) at 13 of 14 stations; they disagree at EC01 (ComPy 0.93, BRUIT-FM 3.51), where one of the two methods failed. Gains from a single event
(Table~\ref{tab:stations}) are less certain than the $\pm$5\,per\,cent prior assumes.
No event passed the coherence test at CC01 and CC09; these stations were kept with the nominal response and a gain sampled over 0.6--1.4, and are excluded from the pooled results.

\begin{figure*}
\centering
\includegraphics[width=0.92\textwidth]{fig05_calibration.png}
\caption{DPG calibration at CC11. (a) Pressure--vertical coherence of the accepted events and its median; dashed: fitted band.
(b) Measured, theoretical and corrected normalized pressure/acceleration ratio.}
\label{fig:cal}
\end{figure*}

\subsection{Compliance measurement}
Compliance is computed on non-overlapping 35-min windows. A window is kept if, in the compliance band, its median coherence is
at least 0.9 (0.8 if fewer than 500 windows pass), the mean calibrated pressure PSD is 10--50\,dB re 1\,Pa$^2$\,Hz$^{-1}$, and the
mean vertical acceleration PSD is between $-170$ and $-150$\,dB re 1\,(m\,s$^{-2}$)$^2$\,Hz$^{-1}$ (Aminian et al.\ 2025).
The band is the widest run of frequencies where the mean coherence of the kept windows reaches the threshold, below the
infragravity cut-off $\sqrt{g/2\pi H}$. Spectra of the kept windows are stacked before the ratio is formed (eq.~1).
Fig.~\ref{fig:coh} shows one month at CC11; the coherence is high whenever the vertical noise is low.
Fig.~\ref{fig:windows} shows the selected windows and the stacked compliance. Monthly stacks with the same band test the
temporal stability.

\begin{figure*}
\centering
\includegraphics[width=0.95\textwidth]{fig06_coherogram.png}
\caption{One month at CC11: spectrograms of calibrated pressure (a) and corrected vertical acceleration (b) for 35-min
windows, coherogram (c), and their band averages (right). Dashed lines: compliance band; shading: windows above the coherence threshold.}
\label{fig:coh}
\end{figure*}

\begin{figure*}
\centering
\includegraphics[width=0.85\textwidth]{fig07_windows.png}
\caption{Selected windows at CC11 (grey, a random subset): (a) vertical acceleration, (b) pressure, (c) coherence,
(d) compliance, with the stacked estimate (red; error bars: window scatter). Dashed: band and coherence threshold.}
\label{fig:windows}
\end{figure*}

\subsection{Shear-velocity inversion}
The inversion follows the parameterization of Zha \& Webb (2016) as implemented in ComPy 2.3. A sediment layer of free thickness
(0--1000\,m) has Hamilton (1976, 1979) Vs and Vp profiles. Below it, crustal layers 100, 200, 400, 800 and 1600\,m thick reach
the CRUST1.0 Moho at the station, above a mantle half-space. Crustal Vs is uniform on 0.5--4.7\,km\,s$^{-1}$ with a hierarchical
second-difference smoothing prior. Crustal Vp/Vs is free on 1.7--2.4 with density from Christensen \& Shaw (1970), and the
pressure gain is sampled within $\pm$5\,per\,cent of the Rayleigh value. Eight chains of delayed-rejection adaptive Metropolis
(Haario et al.\ 2006), each a parallel-tempering ladder of eight temperatures, run 80\,000 iterations (30\,000 burn-in) from
independent prior draws; a prior-only run shows what the data add. Convergence is assessed with the rank-normalized $\hat R$
and effective sample size (Vehtari et al.\ 2021). The variant with Vp/Vs fixed at 1.85 brackets the effect of that
assumption. This differs from the published workflow, whose code is audited in \texttt{ComPy/docs/INVERSION.md}.

\section{Results}
\subsection{Compliance functions}
Table~\ref{tab:stations} and Fig.~\ref{fig:compliance} summarize the measurements. At the 18 calibrated stations the band
is 4.9--14.0\,mHz, the median coherence 0.902--0.970, and the random error about 0.1\,per\,cent.
The functions are remarkably similar across the array: $\eta$ at the lowest frequency is 0.016--0.028\,GPa$^{-1}$ and rises
monotonically with frequency, the signature of Vs increasing with depth. Monthly stacks agree with the full stack to within
1--3\,per\,cent at most stations; CC07 is the exception (-25 to +9\,per\,cent), with a drift that the
Rayleigh-wave gain does not show and which therefore enters its covariance as non-stationarity.

\begin{figure*}
\centering
\includegraphics[width=\textwidth]{fig08_compliance.png}
\caption{Normalized compliance by sub-array, west to east; shading: $\pm1\sigma$ from the full covariance.}
\label{fig:compliance}
\end{figure*}

\begin{table*}
\centering\small
\caption{Stations with a compliance measurement. Gain: ComPy Rayleigh-wave calibration (event standard deviation, number of events); BRUIT-FM: independent microseism calibration where available. Frequency: compliance band. Threshold: coherence gate. Coherence: band median of the stack.}
\label{tab:stations}
\begin{tabular}{lrrrlrrrrr}
\toprule
Station & Lat. [$^\circ$] & Lon. [$^\circ$] & Depth [m] & Gain (ComPy) & Gain (BRUIT-FM) & Band [mHz] & Thr. & Windows & Coh. \\
\midrule
CC01 & -4.0458 & -132.4745 & 4768 & -- & -- & 4.9--13.4 & 0.9 & 2781 & 0.968 \\
CC03 & -4.8816 & -132.6890 & 4739 & 0.763 (1) & 0.913 & 4.9--15.9 & 0.9 & 1149 & 0.965 \\
CC04 & -5.1937 & -132.7900 & 4615 & 0.804 $\pm$ 0.010 (7) & 0.803 & 4.9--14.6 & 0.9 & 7024 & 0.967 \\
CC05 & -5.5174 & -132.8730 & 4798 & 1.060 $\pm$ 0.016 (10) & 1.146 & 4.9--13.4 & 0.9 & 7602 & 0.960 \\
CC06 & -5.8385 & -132.9518 & 4824 & 0.939 $\pm$ 0.019 (6) & 0.965 & 4.9--13.4 & 0.9 & 5002 & 0.964 \\
CC07 & -6.1588 & -133.0464 & 4628 & 0.912 $\pm$ 0.013 (9) & -- & 4.9--13.4 & 0.9 & 5773 & 0.942 \\
CC08 & -6.4816 & -133.1233 & 4745 & 0.839 $\pm$ 0.001 (5) & -- & 4.9--14.6 & 0.9 & 3936 & 0.964 \\
CC09 & -6.7995 & -133.2164 & 4615 & -- & -- & 4.9--17.1 & 0.8 & 501 & 0.957 \\
CC11 & -7.5368 & -133.4166 & 4636 & 0.698 $\pm$ 0.012 (12) & 0.702 & 4.9--14.6 & 0.9 & 7314 & 0.970 \\
EC01 & -4.7261 & -131.4852 & 4576 & 0.927 (1) & 3.509 & 6.1--13.4 & 0.8 & 545 & 0.902 \\
EC05 & -7.8562 & -132.3298 & 4587 & 0.840 $\pm$ 0.012 (5) & 0.824 & 4.9--13.4 & 0.9 & 4093 & 0.961 \\
EE01 & -4.9470 & -130.3814 & 4645 & 0.872 $\pm$ 0.014 (10) & 0.893 & 8.5--13.4 & 0.9 & 3207 & 0.926 \\
EE04 & -8.2043 & -131.2551 & 4683 & 1.146 $\pm$ 0.015 (12) & 1.159 & 4.9--13.4 & 0.9 & 7316 & 0.965 \\
WC02 & -4.9726 & -133.7054 & 4619 & 1.242 $\pm$ 0.009 (6) & 1.259 & 4.9--14.6 & 0.9 & 5074 & 0.967 \\
WC03 & -5.7077 & -134.0913 & 4483 & 0.779 $\pm$ 0.008 (9) & 0.778 & 4.9--14.6 & 0.9 & 8572 & 0.968 \\
WC04 & -6.5352 & -134.1305 & 4534 & 0.865 $\pm$ 0.023 (6) & 0.840 & 4.9--14.6 & 0.9 & 7145 & 0.965 \\
WC05 & -7.2694 & -134.5177 & 4478 & 0.898 $\pm$ 0.015 (9) & 0.992 & 4.9--14.6 & 0.9 & 7222 & 0.961 \\
WW01 & -3.7852 & -134.7370 & 4534 & 0.821 $\pm$ 0.026 (4) & 0.925 & 6.1--12.2 & 0.9 & 1796 & 0.945 \\
WW03 & -6.0155 & -135.1120 & 4386 & 0.873 $\pm$ 0.011 (12) & -- & 6.1--13.4 & 0.9 & 5958 & 0.959 \\
WW04 & -7.0385 & -135.6272 & 4491 & 0.911 $\pm$ 0.003 (3) & -- & 4.9--15.9 & 0.9 & 2736 & 0.964 \\
\bottomrule
\end{tabular}
\end{table*}

\subsection{Shear-velocity profiles}
All inversions converged ($\hat R\leq$ 1.015) and fit to $\chi^2$ below the number of data. Fig.~\ref{fig:heat} shows the
posteriors and Table~\ref{tab:vs} the values at fixed depths. The sediment is 79--210\,m thick, in line with CRUST1.0
(70--190\,m). Below it, Vs rises steeply to 2.35--2.86\,km\,s$^{-1}$ at 0.5\,km and 2.67--3.20\,km\,s$^{-1}$
at 1\,km, then more slowly to 3.01--3.58\,km\,s$^{-1}$ at 2--3\,km. The posterior is 0.36 times as wide as the prior at 1\,km and
0.21 at 2\,km, so these depths are resolved. Half of the sensitivity of the lowest frequency lies above 3.5\,km and
90\,per\,cent above 12.3\,km. Below about 3\,km the posterior width reflects the layer geometry as much as the data.

\begin{figure*}
\centering
\includegraphics[width=\textwidth]{fig09_heatmaps.png}
\caption{Vs posterior heat maps (rows normalized). Red: median; white dashed: CRUST1.0 at the station. Depths are true depths below the seafloor.}
\label{fig:heat}
\end{figure*}

\begin{table*}
\centering\small
\caption{Inversion results: fit ($\chi^2/N$), largest $\hat R$, sediment thickness and crustal Vp/Vs (posterior medians), Vs [km\,s$^{-1}$] as median (5--95 per cent) at depths below the seafloor, and Vs at 4\,km with Vp/Vs = 1.85. $^*$Gain not calibrated (sampled over 0.6--1.4).}
\label{tab:vs}
\resizebox{\textwidth}{!}{\begin{tabular}{lrrrrllllr}
\toprule
Station & $\chi^2/N$ & $\hat R$ & Sed. [m] & Vp/Vs & Vs 0.5\,km & Vs 1\,km & Vs 2\,km & Vs 4\,km & Vs 4\,km (1.85) \\
\midrule
CC01$^*$ & 2.9/8 & 1.015 & 147 & 2.16 & 2.49 (1.79--3.70) & 2.72 (2.14--3.48) & 2.93 (2.51--3.42) & 3.21 (2.74--3.80) & 3.98 \\
CC03 & 3.6/10 & 1.011 & 166 & 2.24 & 2.64 (2.10--3.38) & 2.94 (2.61--3.30) & 3.25 (3.04--3.49) & 3.55 (3.34--3.92) & 4.22 \\
CC04 & 3.1/9 & 1.006 & 81 & 2.23 & 2.64 (2.11--3.48) & 2.85 (2.46--3.33) & 3.09 (2.86--3.33) & 3.35 (3.19--3.63) & 3.89 \\
CC05 & 3.0/8 & 1.006 & 79 & 2.21 & 2.80 (2.18--3.85) & 3.02 (2.52--3.71) & 3.28 (2.93--3.62) & 3.56 (3.37--3.89) & 4.09 \\
CC06 & 3.0/8 & 1.007 & 118 & 2.26 & 2.58 (1.98--3.71) & 2.88 (2.38--3.59) & 3.22 (2.84--3.56) & 3.58 (3.39--3.95) & 4.31 \\
CC07 & 2.8/8 & 1.009 & 137 & 2.17 & 2.83 (2.06--4.09) & 3.20 (2.48--4.02) & 3.58 (3.07--4.10) & 3.99 (3.60--4.55) & 4.44 \\
CC08 & 3.2/9 & 1.015 & 153 & 2.23 & 2.35 (1.79--3.40) & 2.70 (2.23--3.23) & 3.08 (2.78--3.39) & 3.48 (3.24--3.91) & 4.16 \\
CC09$^*$ & 5.4/11 & 1.010 & 85 & 2.12 & 2.27 (1.81--3.01) & 2.64 (2.20--3.12) & 3.08 (2.66--3.49) & 3.52 (3.02--4.07) & 4.17 \\
CC11 & 3.2/9 & 1.009 & 121 & 2.17 & 2.53 (1.94--3.43) & 2.85 (2.39--3.37) & 3.17 (2.89--3.48) & 3.50 (3.26--3.98) & 3.98 \\
EC01 & 3.0/7 & 1.006 & 165 & 2.12 & 2.63 (1.90--3.99) & 2.93 (2.25--3.72) & 3.22 (2.70--3.68) & 3.52 (3.16--4.19) & 3.89 \\
EC05 & 2.9/8 & 1.010 & 133 & 2.20 & 2.61 (1.94--3.74) & 2.90 (2.35--3.58) & 3.24 (2.86--3.58) & 3.58 (3.34--4.01) & 4.18 \\
EE01 & 3.0/5 & 1.010 & 138 & 2.07 & 2.63 (1.92--3.94) & 3.00 (2.29--3.83) & 3.36 (2.78--3.87) & 3.78 (3.39--4.41) & 4.11 \\
EE04 & 2.8/8 & 1.006 & 108 & 2.18 & 2.86 (2.12--3.90) & 3.06 (2.49--3.76) & 3.28 (2.93--3.64) & 3.48 (3.27--3.86) & 3.96 \\
WC02 & 3.0/9 & 1.006 & 112 & 2.16 & 2.72 (2.09--3.65) & 2.96 (2.48--3.53) & 3.20 (2.92--3.48) & 3.43 (3.21--3.83) & 3.88 \\
WC03 & 3.4/9 & 1.005 & 95 & 2.12 & 2.77 (2.12--3.72) & 3.04 (2.55--3.59) & 3.31 (3.01--3.62) & 3.58 (3.33--4.02) & 3.96 \\
WC04 & 3.7/9 & 1.009 & 100 & 2.24 & 2.39 (1.90--3.39) & 2.67 (2.24--3.23) & 3.01 (2.73--3.29) & 3.39 (3.19--3.74) & 4.06 \\
WC05 & 3.1/9 & 1.006 & 139 & 2.16 & 2.61 (2.00--3.51) & 2.90 (2.46--3.40) & 3.18 (2.91--3.47) & 3.46 (3.22--3.92) & 3.94 \\
WW01 & 2.5/6 & 1.009 & 179 & 2.09 & 2.76 (1.96--4.13) & 3.08 (2.37--3.89) & 3.34 (2.81--3.88) & 3.60 (3.24--4.23) & 3.94 \\
WW03 & 2.7/7 & 1.009 & 170 & 2.11 & 2.58 (1.89--3.89) & 2.91 (2.32--3.65) & 3.23 (2.82--3.67) & 3.53 (3.20--4.17) & 3.92 \\
WW04 & 3.7/10 & 1.006 & 210 & 2.23 & 2.76 (2.15--3.84) & 3.13 (2.69--3.63) & 3.35 (3.10--3.66) & 3.57 (3.33--4.02) & 4.19 \\
\bottomrule
\end{tabular}}
\end{table*}

\subsection{Vp/Vs and the lower crust}
The crustal Vp/Vs posteriors lean towards high values (medians 2.07--2.26 on a 1.7--2.4 prior) but remain wide.
Compliance depends on both Lam\'e parameters, so a higher Vp/Vs is traded for a lower Vs. With Vp/Vs fixed at 1.85 the
upper 2\,km barely change, while Vs at 4\,km rises from 3.54 to 4.05\,km\,s$^{-1}$ (network median), close to
CRUST1.0 (4.05\,km\,s$^{-1}$). Both models fit. Lower-crustal Vs is therefore set by the Vp/Vs assumption,
whereas upper-crustal Vs is a robust result.

\subsection{Serpentinization}
Applied to the pooled posterior, eq.~(4) gives 78 (20--100) at 0.5\,km, 63 (27--91) at 1\,km and 49 (25--67)\,per\,cent at
2\,km below the seafloor (median and 5--95\,per\,cent range; Fig.~\ref{fig:prof}c).

\begin{figure*}
\centering
\includegraphics[width=\textwidth]{fig10_profiles.png}
\caption{(a) Median Vs by sub-array (thin lines: stations) and the mean CRUST1.0 profile. (b) Pooled posterior of the calibrated
stations with Vp/Vs free (blue) and the median with Vp/Vs = 1.85 (orange). (c) Serpentinization from eq.~(4) applied to the
pooled draws: the value that \emph{would} follow if the crust were peridotite.}
\label{fig:prof}
\end{figure*}

\section{Discussion}
\textit{A consistent normal-crust signal.} The Young ORCA compliance is uniform across a 600\,km by 500\,km array, and the Vs
profiles agree between stations to within their uncertainties. This is expected from crust formed at a fast-spreading ridge,
where crustal thickness and layering vary little (e.g.\ Christeson et al.\ 2019). The velocity structure is the classical one:
a thin pelagic sediment cover; a steep gradient through the extrusive section (layer 2A--2B), where Vs climbs from
$\sim$2.3 to $\sim$3\,km\,s$^{-1}$ in the first kilometre as porosity and crack density decrease (Carlson 2014); and a gentler
gradient in the dikes and gabbros. Upper-crustal Vs is lower than in CRUST1.0, which averages refraction models and smooths
the steep gradient at the top of the crust.

\textit{What this means for serpentinization.} Fed through the Miller \& Christensen (1997) relation, the upper-crustal
velocities at Young ORCA would read as 63\,per\,cent serpentinization at 1\,km, comparable with the values now
obtained beneath the SWIR stations of RHUM-RUM (30--55\,per\,cent at 2\,km with the corrected relation and models;
ComPy 2.4 report). Serpentinite is not expected in this magmatic crust; the low Vs is the porosity of the extrusives.
Vs from compliance alone therefore cannot distinguish serpentinized peridotite from porous basalt. The SWIR interpretation
has to rest on the geological setting (exhumed mantle in oblique, magma-poor sections; Sauter \& Cannat 2010) and, where possible,
on independent constraints: Vp/Vs, which is high in serpentinite ($\geq$2.0; Christensen 2004) but is poorly resolved by
compliance (Section~5.3), density, or active-source data. A useful diagnostic is the \emph{difference} between ridge and normal
crust at equal depth. At 2\,km below the seafloor, Young ORCA gives 3.23\,km\,s$^{-1}$; the corrected models of RR36 and
RR40 give 3.09 and 3.56\,km\,s$^{-1}$ and RR52 (CIR) 3.59\,km\,s$^{-1}$. None of them is slower than the Young ORCA range (3.01--3.58\,km\,s$^{-1}$ at 2\,km), so at that depth the ridge stations are not distinguishable from normal crust by Vs alone.

\textit{Limitations.} (i) The upper bound of the Vp/Vs prior is reached by several posteriors, so a wider prior could lower
lower-crustal Vs further. (ii) The compliance of CC07 drifts slowly; the drift widens its covariance and posterior. CC01 and CC09 have no
Rayleigh-wave calibration (gain sampled over 0.6--1.4); their profiles are shown but not pooled. EC01 and CC03 rest on a single
calibration event. (iii) A 1-D model ignores the abyssal-hill topography, partly covered by the
1\,per\,cent covariance floor (Hulme et al.\ 2005). (iv) Stations with long clipped spans have short records.

\section{Conclusions}
ComPy 2.4.0 processed a second broad-band OBS network from the raw archive to compliance, shear velocity and
serpentinization without manual intervention, after adding the removal of clipped data that this network required. At
20 Young ORCA stations it measures compliance with coherence above 0.9 over 4.9--14.0\,mHz, and its
Rayleigh-wave DPG calibration agrees with an independent microseism calibration. The inverted crust of $\sim$40\,Ma Pacific
lithosphere has 79--210\,m of sediment and upper-crustal Vs of 2.35--2.86\,km\,s$^{-1}$ at 0.5\,km. Lower-crustal Vs depends
on Vp/Vs. These normal-crust velocities would translate into substantial ``serpentinization'' if read like the SWIR results, so
compliance-derived Vs must be combined with geological and independent geophysical constraints before it is interpreted as alteration.

\section*{Data and software}
Young ORCA data (FDSN network XE, 2018--2019) are distributed by the EarthScope Consortium (IRIS DMC). ComPy
2.4.0, including the pipeline and the scripts of this run, is at \url{https://github.com/MohammadAmin-Aminian/ComPy}.
Run folder: \texttt{BRUIT-FM/compliance\_xe\_2026-10-29}.

\footnotesize
\section*{References}
\begin{description}\setlength{\itemsep}{0pt}
\item Aminian, M.A., Crawford, W., Stutzmann, \'E., Montagner, J.-P., Cannat, M. \& Hadziioannou, C., 2025. Shallow crustal structures of the Indian ocean derived from compliance function analysis, \textit{Geophys. J. Int.}, 242, ggaf253.
\item Bendat, J.S. \& Piersol, A.G., 1993. \textit{Engineering Applications of Correlation and Spectral Analysis}, Wiley.
\item Carlson, R.L., 2014. The effects of alteration and porosity on seismic velocities in oceanic basalts and diabases, \textit{Geochem. Geophys. Geosyst.}, 15.
\item Christensen, N.I., 2004. Serpentinites, peridotites, and seismology, \textit{Int. Geol. Rev.}, 46, 795--816.
\item Christensen, N.I. \& Shaw, G.H., 1970. Elasticity of mafic rocks from the Mid-Atlantic Ridge, \textit{Geophys. J. R. astr. Soc.}, 20, 271--284.
\item Christeson, G.L., Goff, J.A. \& Reece, R.S., 2019. Synthesis of oceanic crustal structure from two-dimensional seismic profiles, \textit{Rev. Geophys.}, 57, 504--529.
\item Crawford, W.C. \& Webb, S.C., 2000. Identifying and removing tilt noise from low-frequency ($<$0.1\,Hz) seafloor vertical seismic data, \textit{Bull. seism. Soc. Am.}, 90, 952--963.
\item Crawford, W.C., Webb, S.C. \& Hildebrand, J.A., 1991. Seafloor compliance observed by long-period pressure and displacement measurements, \textit{J. geophys. Res.}, 96, 16\,151--16\,160.
\item Crawford, W.C., Webb, S.C. \& Hildebrand, J.A., 1998. Estimating shear velocities in the oceanic crust from compliance measurements by two-dimensional finite difference modeling, \textit{J. geophys. Res.}, 103, 9895--9916.
\item Deen, M. et al., 2017. First observation of the Earth's permanent free oscillations on ocean bottom seismometers, \textit{Geophys. Res. Lett.}, 44, 10\,988--10\,996.
\item Eilon, Z.C. et al., 2022. The Pacific OBS Research into Convecting Asthenosphere (ORCA) experiment, \textit{Seismol. Res. Lett.}
\item Gaherty, J., Eilon, Z., Forsyth, D. \& Ekstr\"om, G., 2019. PacificORCA, International Federation of Digital Seismograph Networks, network XE.
\item Haario, H., Laine, M., Mira, A. \& Saksman, E., 2006. DRAM: efficient adaptive MCMC, \textit{Stat. Comput.}, 16, 339--354.
\item Hamilton, E.L., 1976. Shear-wave velocity versus depth in marine sediments: a review, \textit{Geophysics}, 41, 985--996.
\item Harmon, N., Laske, G., Crawford, W. \& Rychert, C., 2022. Tilt corrections for normal mode observations on ocean bottom seismic data, \textit{Seismica}, 1.
\item Hulme, T., Crawford, W.C. \& Singh, S.C., 2005. The sensitivity of seafloor compliance to two-dimensional low-velocity anomalies, \textit{Geophys. J. Int.}, 163, 547--558.
\item Laske, G., Masters, G., Ma, Z. \& Pasyanos, M., 2013. Update on CRUST1.0, \textit{EGU General Assembly}, 15, 2658.
\item Miller, D.J. \& Christensen, N.I., 1997. Seismic velocities of lower crustal and upper mantle rocks from the slow-spreading Mid-Atlantic Ridge, \textit{Proc. ODP Sci. Results}, 153, 437--454.
\item Sauter, D. \& Cannat, M., 2010. The ultraslow spreading Southwest Indian Ridge, \textit{Geophys. Monogr.}, 188, AGU.
\item Sorrells, G.G. \& Goforth, T.T., 1973. Low-frequency earth motion generated by slowly propagating partially organized pressure fields, \textit{Bull. seism. Soc. Am.}, 63, 1583--1601.
\item Vehtari, A., Gelman, A., Simpson, D., Carpenter, B. \& B\"urkner, P.-C., 2021. Rank-normalization, folding, and localization: an improved $\hat R$, \textit{Bayesian Anal.}, 16, 667--718.
\item Zha, Y. \& Webb, S.C., 2016. Crustal shear velocity structure in the Southern Lau Basin constrained by seafloor compliance, \textit{J. geophys. Res.}, 121, 3220--3237.
\end{description}
\end{document}
